\section{Conclusion}
\label{sec:conclusion}

% Included in both manuscript versions on 2026-11-03 at the author's request.
In this paper, we investigated GNSS-independent GFRA for LEO satellite
IoT systems. We proposed a multi-frame AFDM transmission scheme that
accommodates large differential delays and Doppler shifts without
device-side precompensation. Under a shared effective-channel
representation, additional pilot frames provide more observations without
increasing the number of channel coefficients to be estimated. We learned
a score-based generative prior for AFDM effective channels without imposing
sparse or grid-aligned channel assumptions, and developed a denoiser that
provides channel estimates and average variance without requiring a
separately trained variance model. We incorporated this denoiser into a
message-passing JADCE receiver comprising MFLE and GPJD. The learned
density-ratio estimator in GPJD enables activity detection under the
learned prior without explicit marginal-likelihood evaluation. Simulation
results demonstrated improved activity detection and channel estimation
over the adapted STMP baseline, with fewer pilot frames required for
comparable channel-estimation performance. The receiver also achieved
lower NMSE than the sparse-recovery baselines across the tested NTN-TDL
profiles without retraining.
